%------------------------
% Resume in Latex
% Generated from content/resume.zh.json
%------------------------

\documentclass[letterpaper,11pt]{article}

\usepackage{latexsym}
\usepackage[empty]{fullpage}
\usepackage{titlesec}
\usepackage{marvosym}
\usepackage[usenames,dvipsnames]{color}
\usepackage{verbatim}
\usepackage{enumitem}
\usepackage[hidelinks]{hyperref}
\usepackage{fancyhdr}
\usepackage[english]{babel}
\usepackage{tabularx}
\usepackage[UTF8,scheme=plain,fontset=none]{ctex}
\setCJKmainfont{AR PL UMing TW}
\setCJKsansfont{AR PL UMing TW}
\setCJKmonofont{AR PL UMing TW}
\usepackage{fontawesome5}
\usepackage[scale=0.90,lf]{FiraMono}
\usepackage{contour}
\usepackage[normalem]{ulem}
\usepackage{tgheros}
\usepackage[T1]{fontenc}

\definecolor{light-grey}{gray}{0.83}
\definecolor{dark-grey}{gray}{0.3}
\definecolor{text-grey}{gray}{.08}

\DeclareRobustCommand{\ebseries}{\fontseries{eb}\selectfont}
\DeclareTextFontCommand{\texteb}{\ebseries}

\renewcommand{\ULdepth}{1.8pt}
\contourlength{0.8pt}
\newcommand{\myuline}[1]{%
  \uline{\phantom{#1}}%
  \llap{\contour{white}{#1}}%
}

\renewcommand*\familydefault{\sfdefault}

\pagestyle{fancy}
\fancyhf{}
\fancyfoot{}
\renewcommand{\headrulewidth}{0pt}
\renewcommand{\footrulewidth}{0pt}

\addtolength{\oddsidemargin}{-0.5in}
\addtolength{\evensidemargin}{0in}
\addtolength{\textwidth}{1in}
\addtolength{\topmargin}{-.5in}
\addtolength{\textheight}{1.0in}

\urlstyle{same}
\raggedbottom
\raggedright
\setlength{\tabcolsep}{0in}

\titleformat {\section}{
    \bfseries \vspace{2pt} \raggedright \large
}{}{0em}{}[\color{light-grey} {\titlerule[2pt]} \vspace{-4pt}]

\newcommand{\resumeItem}[1]{
  \item\small{
    {#1 \vspace{-1pt}}
  }
}

\newcommand{\resumeSubheading}[4]{
  \vspace{-1pt}\item
    \begin{tabular*}{\textwidth}[t]{l@{\extracolsep{\fill}}r}
      \textbf{#1} & {\color{dark-grey}\small #2}\vspace{1pt}\\
      \textit{#3} & {\color{dark-grey} \small #4}\\
    \end{tabular*}\vspace{-4pt}
}

\newcommand{\resumeProjectHeading}[2]{
    \item
    \begin{tabular*}{\textwidth}{l@{\extracolsep{\fill}}r}
      #1 & {\color{dark-grey} #2} \\
    \end{tabular*}\vspace{-4pt}
}

\renewcommand\labelitemii{$\vcenter{\hbox{\tiny$\bullet$}}$}

\newcommand{\resumeSubHeadingListStart}{\begin{itemize}[leftmargin=0in, label={}]}
\newcommand{\resumeSubHeadingListEnd}{\end{itemize}}
\newcommand{\resumeItemListStart}{\begin{itemize}}
\newcommand{\resumeItemListEnd}{\end{itemize}\vspace{0pt}}

\color{text-grey}
\setlength{\footskip}{5pt}

\begin{document}

\begin{center}
    \textbf{\Huge Joe (永鴻) 黃} \\ \vspace{5pt}
    System / DevOps Engineer \\ \vspace{4pt}
    \hspace{1pt} \faEnvelope \hspace{2pt} \texttt{\href{mailto:yunghunghuang984@gmail.com}{yunghunghuang984@gmail.com}} \hspace{1pt} $|$
    \hspace{1pt} \faGithub \hspace{2pt} \texttt{\href{https://github.com/blackdesert575}{blackdesert575}} \hspace{1pt} $|$
    \hspace{1pt} \faMapMarker* \hspace{2pt}\texttt{Taipei, Taiwan}
    \\ \vspace{-3pt}
\end{center}

\section{摘要}
\begin{itemize}[leftmargin=0in, label={}]
  \item\small{具備 Linux 維運、雲端與地端基礎設施、Kubernetes、CI/CD、觀測性與基礎設施自動化經驗的 System / DevOps / SRE 工程師。}
\end{itemize}

\section{工作經驗}
\resumeSubHeadingListStart
\resumeSubheading
  {System / DevOps Engineer}{2025/06 -- 至今}
  {雷技資訊科技有限公司}{}
  \resumeItemListStart
    \resumeItem{在 Linux 終端機環境中工作。}
    \resumeItem{使用 Packer 構築與維護以 Rocky Linux 作為基底的 EC2 機器映像檔。}
    \resumeItem{主要在台灣工程團隊內維運 AWS 與地端基礎設施，並實際導入與使用以 Git 為核心的 Infrastructure as Code 工作流程，包含 Terraform、Terragrunt 與 Atlantis。}
    \resumeItem{協助審查與驗證既有技術文件與基礎設施流程，並辨識跨區域之間流程與標準不一致的問題。}
    \resumeItem{維護透過 Ansible 進行的系統初始化自動化，包含 Linux 設定、Java 執行環境、Nginx、RabbitMQ 與 NFS 部署。}
    \resumeItem{為台灣工程團隊建立與維護基於 Groovy Script 的 Jenkins CI/CD 流水線。}
    \resumeItem{管理 Bastion 主機與內部資產追蹤系統，基於 Consul 與 Golang 微服務，並將 EC2 metadata 整合至集中式存取控制機制。}
    \resumeItem{透過 Prometheus 與 Grafana 分析系統資源使用情況與效能指標，以識別基礎設施瓶頸並提出優化方案。}
    \resumeItem{維護以 Prometheus、Grafana、Consul 組成的觀測性技術堆疊。}
    \resumeItem{\textbf{技術}: Linux, AWS, Packer, Terraform, Terragrunt, Atlantis, Ansible, Jenkins, Groovy, Prometheus, Grafana, Consul, Golang}
  \resumeItemListEnd

\resumeSubheading
  {System / Site Reliability Engineering (SRE) / DevOps Engineering}{2023/11 -- 2024/12}
  {眾鼎科技有限公司}{}
  \resumeItemListStart
    \resumeItem{在 Linux 終端機環境中工作。}
    \resumeItem{參與輪值待命，約每月一週，與團隊成員共同排除正式環境中的問題。}
    \resumeItem{支援部署於 AWS 或地端基礎設施上的加密貨幣交易所平台。}
    \resumeItem{管理與維運部署於 AWS 或地端環境的 Kubernetes 叢集。}
    \resumeItem{使用 Pulumi 與 Ansible 進行 AWS 與 Cloudflare 基礎設施管理。}
    \resumeItem{管理與維運基於 GitLab CI 或 Jenkins 的 CI/CD 流水線。}
    \resumeItem{\textbf{技術}: Linux, AWS, Cloudflare, Kubernetes, Pulumi, Ansible, GitLab CI, Jenkins}
  \resumeItemListEnd

\resumeSubheading
  {Site Reliability Engineering (SRE) / DevOps / Cloud Engineering}{2021/11 -- 2023/04}
  {雲磐科技有限公司}{}
  \resumeItemListStart
    \resumeItem{在 Linux 終端機環境中工作。}
    \resumeItem{參與輪值待命，與團隊成員共同排除各類系統問題。}
    \resumeItem{開發基於 Python 的網站監控自動化專案。}
    \resumeItem{開發基於 GCP GKE 與阿里雲 ACK 的串流平台專案。}
    \resumeItem{透過 JumpServer 管理 AWS、GCP、阿里雲、Cloudflare 等雲端基礎設施。}
    \resumeItem{管理與維運部署於 AWS EKS、GCP GKE、阿里雲 ACK 的 Kubernetes 叢集。}
    \resumeItem{管理與維運基於 GitLab CI、Jenkins、Tekton 的 CI/CD 流水線。}
    \resumeItem{\textbf{技術}: Linux, Python, GCP, GKE, Aliyun, ACK, AWS, EKS, Cloudflare, JumpServer, Kubernetes, GitLab CI, Jenkins, Tekton}
  \resumeItemListEnd

\resumeSubheading
  {物聯網整合開發人才就業班}{2021/01 -- 2021/05}
  {工業技術研究院 產業學院}{}
  \resumeItemListStart
    \resumeItem{完成一項密集的實作導向培訓課程，內容涵蓋網頁開發基礎、系統程式設計、Linux 環境、網路基礎，以及 IoT 系統整合。}
    \resumeItem{使用 Raspberry Pi 與 ESP8266 建置 IoT 監控系統 IndoorAirBox，專注於溫度與濕度資料蒐集。}
    \resumeItem{\textbf{技術}: Web Development, Linux, Networking, IoT, Raspberry Pi, ESP8266}
  \resumeItemListEnd
\resumeSubHeadingListEnd

\section{專案}

\resumeSubHeadingListStart
\resumeProjectHeading
  {\textbf{homelab}}{}
  \resumeItemListStart
    \resumeItem{建置並維運基於 Proxmox VE 的個人 Linux 實驗環境，以支援多個個人專案。}
    \resumeItem{建置基於 Proxmox VE 的基礎設施，以支援多個個人專案的 Linux Lab 環境。}
    \resumeItem{\textbf{技術}: Proxmox VE, Linux}
  \resumeItemListEnd
\resumeSubHeadingListEnd

\resumeSubHeadingListStart
\resumeProjectHeading
  {\textbf{resume}}{}
  \resumeItemListStart
    \resumeItem{基於 LaTeX 建置的線上履歷網站，用於記錄工作經歷、專案成果、學歷背景與技能。}
    \resumeItem{建置基於 LaTeX 的線上履歷網站。}
    \resumeItem{\textbf{原始碼}: \href{https://github.com/blackdesert575/resume}{https://github.com/blackdesert575/resume}}
    \resumeItem{\textbf{技術}: LaTeX, GitHub Actions, Cloudflare Pages}
  \resumeItemListEnd
\resumeSubHeadingListEnd

\section{學歷}
\resumeSubHeadingListStart
\resumeSubheading
  {國立東華大學}{2016 -- 2019}
  {材料科學與工程學系 碩士}{壽豐鄉, 花蓮縣}

\resumeSubheading
  {國立東華大學}{2012 -- 2016}
  {材料科學與工程學系 學士}{壽豐鄉, 花蓮縣}
\resumeSubHeadingListEnd

%-----------PROGRAMMING SKILLS-----------
\section{技能}
 \begin{itemize}[leftmargin=0in, label={}]
    \small{\item{
     \textbf{語言}{: Mandarin (native), English (professional working proficiency)} \\
     \textbf{程式語言}{: Python (proficient), Bash (competence), SQL (competence), Rust (competence), Golang (beginner)} \\
     \textbf{版本控制}{: Git} \\
     \textbf{軟體框架}{: Django, FastAPI, Selenium, pandas, SQLAlchemy} \\
     \textbf{作業系統}{: Linux (proficient), macOS (competence), Windows (beginner)} \\
     \textbf{持續整合與持續部署}{: GitLab Runner, GitHub Actions, Jenkins, Tekton Pipelines, Argo CD} \\
     \textbf{串流與訊息傳遞}{: Apache RocketMQ, RabbitMQ} \\
     \textbf{容器排程與編排}{: Kubernetes (EKS, GKE, ACK, k0s, k3s)} \\
     \textbf{Kubernetes 子系統}{: ingress-nginx, cert-manager, ECK operator, KEDA, Istio} \\
     \textbf{服務代理}{: NGINX, HAProxy, MetalLB, Envoy} \\
     \textbf{微服務發現與協調}{: Apache Zookeeper, Nacos, Consul} \\
     \textbf{容器映像檔倉庫}{: Harbor} \\
     \textbf{自動化配置}{: Ansible, Atlantis, Terraform, Terragrunt, Pulumi} \\
     \textbf{雲端供應商}{: AWS, Google Cloud, Alibaba Cloud, Cloudflare} \\
     \textbf{專業認證}{: AWS Certified Solutions Architect Associate, AWS Certified Cloud Practitioner}
    }}
 \end{itemize}

\end{document}
